\documentclass[10pt,a4paper]{amsart}
\usepackage[margin=19mm]{geometry}
\usepackage{amsmath,amssymb,mathtools}
\usepackage[scr=rsfs,cal=euler]{mathalfa}
\usepackage{xfrac}
\usepackage[unicode,hidelinks,bookmarks=false]{hyperref}
\newcommand{\ie}{i.e.\ }
\newcommand{\cf}{cf.\ }
\newcommand{\resp}{respectively\ }
\newcommand{\Subsection}{\subsection}
\providecommand{\nobreakdash}{\nobreak\hbox{-}}
\setlength{\emergencystretch}{2em}
\multlinegap0pt
\usepackage{amssymb}
\usepackage{mathtools}
\usepackage{multirow}
\newtheorem{proposition}{Proposition}
\newtheorem{lemma}{Lemma}
\newcommand{\A}{\mathcal{A}}
\newcommand{\N}{\mathbb{N}}
\newcommand{\R}{\mathbb{R}}
\newcommand{\dej}{\delta_j}
\newcommand{\e}{\varepsilon}
\newcommand{\ej}{\varepsilon_j}
\newcommand{\f}{\varphi}
\title{Multiscale homogenization of non-local energies \\ of convolution-type}
\author{Giuseppe Cosma Brusca}
\date{Source: arXiv:2512.18697v1 (CC BY 4.0)}
\begin{document}
\maketitle

\noindent\emph{Source and licence.} Multiscale homogenization of non-local energies \\ of convolution-type. Version arXiv:2512.18697v1 (CC BY 4.0), distributed by arXiv under the Creative Commons Attribution 4.0 International licence, http://creativecommons.org/licenses/by/4.0/, as recorded in the arXiv licence metadata for this exact version and shown on the abstract page https://arxiv.org/abs/2512.18697v1. Full licence notice: \texttt{licenses/arxiv-2512.18697-CC-BY-4.0.txt}.\par\smallskip

\noindent\textit{Editorial scope.} Counted unit: the article's proposition formula (source0.tex lines 608--632). The article's complete original proof of it is delivered with it (source0.tex lines 633--769, one \texttt{proof} environment of 137 lines). No mathematics is written, repaired or re-derived: the statement and the proof are the article's own lines, in the article's own notation, and no retained line is substituted or re-flowed.  Retained uncounted as the source-local context the proof needs: lines 165--168; lines 170--173; lines 175--178; lines 179--182; lines 184--187; lines 462--471; lines 541--553.  Nothing was omitted from any retained span, so the package carries no omission record.  No retained line contains a citation, so the wrapper carries no bibliography and no bibliography entry.  No retained line contains a figure environment, an image-inclusion command or drawing code, so no image is delivered and the package declares no asset dependency.  Editorial formatting only, in addition: an AMS \texttt{amsart} preamble that restates the article's own declarations and macros used by the retained lines, namely the environments \texttt{proposition}, \texttt{lemma} on separate counters, together with the standard packages those lines need.  The retained environments print in this document as Proposition 1, Lemma 1 in order of appearance, whereas the article prints the same items with the numbering of its own full document.  The complete native source of the article as distributed in the arXiv e-print for this exact version is bundled unchanged as \texttt{sources/191412/source0.tex}.
\begin{equation}\label{rho : ball}
\tag{$\rho_1$}
    \rho(\xi)\geq c_0 \quad \text{ for almost every } \xi\in B_{r_0},
\end{equation}

\begin{equation}\label{rho : moment}
\tag{$\rho_2$}
    \int_{\R^d}\rho(\xi)|\xi|^p\, d\xi <+\infty,
\end{equation}

\begin{align}  \notag
& f(\cdot,y,z) \text{ and } f(x,\cdot,z)  \text{ are } Q_1\text{-periodic for every } z\in \R^m \\ \label{f : period} \tag{P}
& \qquad\text{ and for almost every } y, x\in\R^d, \text{ respectively},
\end{align}

\begin{equation} \label{f : convex}
    \tag{C}
f(x,y,\cdot) \text{ is convex for almost every } x,y\in \R^d,
\end{equation}

\begin{equation} \label{f : growth}
    \tag{GC}
\alpha|z|^p\leq f(x,y,z)\leq \beta|z|^p \text{ for almost every } x,y\in \R^d \text{ and for every } z\in \R^m
\end{equation}

\begin{proposition}\label{prop : boundary} Let $T,s>0$, $M\in\R^{m\times d}$, and assume there exists a functional $F: L^p(\Omega;\R^m)\times \A_{\rm reg}(\Omega)\to[0,+\infty]$ such that 
\begin{equation*}
    \Gamma(L^p)\text{-}\lim_{j\to+\infty} F^T_j(u,A) = F(u,A) 
\end{equation*}
    for every $u\in L^p(\Omega; \R^m)$ and $A\in \A_{\rm reg}(\Omega)$. Then
    \begin{equation*}
        \lim_{j\to+\infty}\inf\{F^T_j(u,A) : u\in \mathcal{D}_{s\e_j,M}(A; \R^m)\} = \inf \{F(u,A) : u-Mx\in W^{1,p}_0(A; \R^m)\}
    \end{equation*}
    for every $A\in \A_{\rm reg}(\Omega)$.
\end{proposition}

\begin{lemma} \label{lemma : compactness}
    Let $M\in \R^{m\times d}$, $\{u_j\}_j\subset L^p_{\#,M}(Q_1;\R^m)$, $\{\lambda_j\}_j\subset \R^+$ such that $\lambda_j\to \lambda\in[0,+\infty)$ as $j\to+\infty$, and assume there exists $A$ a bounded open subset of $\R^d$ with Lipschitz boundary such that $Q_1\subseteq A$ and
    \begin{equation*}
        \sup_{j} \Bigl\{ \int_{Q_1} u_j\, dx + \int_{B_r}\int_{A}\Bigl|\frac{u_j(x+\lambda_j\xi)-u_j(x)}{\lambda_j}\Bigr|^p\, dx\, d\xi \Bigr\}<+\infty
    \end{equation*}
    for some $r>0$.
    The following hold:
    \begin{itemize}
        \item[(i)] if $\lambda=0$, there exist a subsequence $\{u_{j_k}\}_k$ and a function $u\in W^{1,p}(A;\R^m)$ such that $u_{j_k}\to u$ strongly in $L^p(A;\R^m)$ as $k\to+\infty$;
        \item[(ii)] if $\lambda\in(0,+\infty)$, there exist a subsequence $\{u_{j_k}\}_k$ and a function $u\in L^{p}(A;\R^m)$ such that $u_{j_k}\rightharpoonup u$ weakly in $L^p(A;\R^m)$ as $k\to+\infty$;
    \end{itemize}
    in both cases, it holds that $\sup_j \|u_j\|_{L^p(Q_1;\R^m)}$ is finite.
\end{lemma}

\begin{proposition}\label{prop : cell formula}
    Assume there exists a quasiconvex function $\varphi: \R^{m\times d}\to[0,+\infty)$ such that 
\begin{equation*}
     \Gamma(L^p)\text{-}\lim_{j\to+\infty} F^T_j(u,A) = \begin{cases}
            \displaystyle \int_{A} \f(\nabla u)\,dx & \text{ if } u\in W^{1,p}(A;\R^m), \\
            +\infty & \text{ if } u\in L^p(\Omega;\R^m)\setminus W^{1,p}(A;\R^m)
        \end{cases}
\end{equation*}
for every $A\in\A_{\rm reg}(\Omega)$. For every $M\in\R^{m\times d}$ the following hold:
\begin{itemize}
    \item[(i)]  for every $\lambda\in[0,+\infty]$ we have
\end{itemize}\begin{align*}
    \f(M) & = \lim_{j\to+\infty} \inf\Bigl\{ \Bigl(\frac{\dej}{r}\Bigr)^d\int_{B_T} \rho(\xi)\int_{(Q_{r/\dej})_{\lambda_j}(\xi)} f\Bigl(x, x+\lambda_j\xi, \frac{u(x+\lambda_j\xi)-u(x)}{\lambda_j}\Bigr)\, dx\,d\xi : \\
    & \qquad \qquad \qquad u\in \mathcal{D}_{T\lambda_j,M}(Q_{r/\dej};\R^m) \Bigr\}
\end{align*}
for some $r>0$;
\item[(ii)] for every $\lambda\in[0,+\infty]$ we have 
\begin{equation*}
    \f(M) \geq \limsup_{j\to+\infty} \inf\Bigl\{ \int_{B_T} \rho(\xi)\int_{Q_1} f\Bigl(x,x+\lambda_j\xi,\frac{u(x+\lambda_j\xi)-u(x)}{\lambda_j}\Bigr)\, dx\,d\xi : u\in L^p_{\#,M}(Q_1;\R^m) \Bigr\};
\end{equation*}
\item[(iii)] for every $\lambda\in[0,+\infty)$ we have
\begin{equation*}
    \f(M) = \lim_{j\to+\infty} \inf\Bigl\{ \int_{B_T} \rho(\xi)\int_{Q_1} f\Bigl(x,x+\lambda_j\xi,\frac{u(x+\lambda_j\xi)-u(x)}{\lambda_j}\Bigr)\, dx\,d\xi : u\in L^p_{\#,M}(Q_1;\R^m) \Bigr\}.
\end{equation*}
\end{proposition}

\begin{proof}
Fix $M\in \R^{m\times d}$ and let $Q$ be a cube of side-length $r$ which is contained in $\Omega$. Using the quasiconvexity of $\f$ and the convergence of boundary-value problems established in Proposition \ref{prop : boundary} with $A=Q$ and $s=T$, we have
    \begin{align*}
      \f(M) & = \inf \Bigl\{ \frac{1}{r^d}\int_{Q} \f(\nabla u) \, dx: u-Mx\in W^{1,p}_0(Q; \R^m) \Bigr\} \\ 
      & = \lim_{j\to+\infty} \inf\Bigl\{ \frac{1}{r^d}F^T_j(u, Q) : u\in \mathcal{D}_{T\ej,M}(Q; \R^m) \Bigr\}.
    \end{align*}
For the sake of exposition, we assume that $Q=Q_r=(0,r)^d$, the general case being analogous. Consider $u\in \mathcal{D}_{T\ej,M}(Q_r;\R^m)$, and set $v(x):=u(\dej x)/\dej$. We have that $v\in \mathcal{D}_{T\lambda_j,M}(Q_{r/\dej}; \R^m)$ and, by a change of variables, 
\begin{align*}
    F_j^T(u, Q_r) & = \int_{B_T} \rho(\xi)\int_{(Q_r)_{\ej}(\xi)} f\Bigl(\frac{x}{\dej},\frac{x+\ej\xi}{\dej},\frac{u(x+\ej\xi)-u(x)}{\ej}\Bigr)\, dx\,d\xi \\
    & = \int_{B_T} \rho(\xi)\int_{(Q_r)_{\ej}(\xi)} f\Bigl(\frac{x}{\dej},\frac{x+\ej\xi}{\dej},\frac{v(\tfrac{x}{\dej}+\tfrac{\ej}{\dej}\xi)-v(\tfrac{x}{\dej})}{\ej/\dej}\Bigr)\, dx\,d\xi \\
    & = \dej^d\int_{B_T} \rho(\xi)\int_{(Q_{r/\dej})_{\lambda_j}(\xi)} f\Bigl(x, x+\lambda_j\xi, \frac{v(x+\lambda_j\xi)-v(x)}{\lambda_j}\Bigr)\, dx\,d\xi.
\end{align*}
By the arbitrariness of $u$, we obtain that 
\begin{align*}
    \f(M) & = \lim_{j\to+\infty} \inf\Bigl\{ \Bigl(\frac{\dej}{r}\Bigr)^d\int_{B_T} \rho(\xi)\int_{(Q_{r/\dej})_{\lambda_j}(\xi)} f\Bigl(x, x+\lambda_j\xi, \frac{u(x+\lambda_j\xi)-u(x)}{\lambda_j}\Bigr)\, dx\,d\xi : \\
    & \qquad \qquad \qquad u\in \mathcal{D}_{T\lambda_j,M}(Q_{r/\dej}; \R^m) \Bigr\},
\end{align*}
which proves $(i)$.

\smallskip

Now we let
\begin{equation*}
    \f_j(M):=\inf\Bigl\{ \int_{B_T} \rho(\xi)\int_{Q_1} f\Bigl(x,x+\lambda_j\xi,\frac{u(x+\lambda_j\xi)-u(x)}{\lambda_j}\Bigr)\, dx\,d\xi : u\in L^p_{\#,M}(Q_1; \R^m) \Bigr\}
\end{equation*}
and prove that
\begin{equation}\label{cell lemma}
    \f(M)\geq \limsup_{j\to+\infty} \f_j(M).
\end{equation}

Given $u\in \mathcal{D}_{T\lambda_j,M}(Q_{r/\dej}; \R^m)$, we define
\begin{equation*}
    \widetilde{u}(x):=
    \begin{cases}
        u(x) & \text{ if } x\in Q_{r/\dej}, \\
        Mx & \text{ if } x\in Q_{\lceil r/\dej\rceil}\setminus Q_{r/\dej},
    \end{cases}
\end{equation*}
we note that $\widetilde{u}\in \mathcal{D}_{T\lambda_j,M}(Q_{\lceil r/\dej\rceil}; \R^m)$, and we write
\begin{align*}
    &\Bigl(\frac{\dej}{r}\Bigr)^d\int_{B_T} \rho(\xi)\int_{(Q_{\lceil r/\dej \rceil })_{\lambda_j}(\xi)} f\Bigl(x, x+\lambda_j\xi, \frac{\widetilde{u}(x+\lambda_j\xi)-\widetilde{u}(x)}{\lambda_j}\Bigr)\, dx\,d\xi \\
    & = \Bigl(\frac{\dej}{r}\Bigr)^d\int_{B_T} \rho(\xi)\int_{(Q_{r/\dej})_{\lambda_j}(\xi)} f\Bigl(x, x+\lambda_j\xi, \frac{u(x+\lambda_j\xi)-u(x)}{\lambda_j}\Bigr)\, dx\,d\xi \\
    & \quad + \Bigl(\frac{\dej}{r}\Bigr)^d\int_{B_T} \rho(\xi)\int_{(Q_{\lceil r/\dej \rceil })_{\lambda_j}(\xi) \setminus (Q_{r/\dej})_{\lambda_j}(\xi)} f\Bigl(x, x+\lambda_j\xi, \frac{\widetilde{u}(x+\lambda_j\xi)-\widetilde{u}(x)}{\lambda_j}\Bigr)\, dx\,d\xi.
\end{align*}
If $\xi\in B_T$ and $x\in (Q_{\lceil r/\dej \rceil })_{\lambda_j}(\xi) \setminus (Q_{r/\dej})_{\lambda_j}(\xi)$, then $\widetilde{u}(x+\lambda_j\xi)-\widetilde{u}(x)=\lambda_j M\xi$, hence, using the upper bound in \eqref{f : growth} we get
\begin{align} \nonumber
 & \Bigl(\frac{\dej}{r}\Bigr)^d\int_{B_T} \rho(\xi)\int_{(Q_{\lceil r/\dej \rceil })_{\lambda_j}(\xi) \setminus (Q_{r/\dej})_{\lambda_j}(\xi)} f\Bigl(x, x+\lambda_j\xi, \frac{\widetilde{u}(x+\lambda_j\xi)-\widetilde{u}(x)}{\lambda_j}\Bigr)\, dx\,d\xi \\ \nonumber
    & = \Bigl(\frac{\dej}{r}\Bigr)^d\int_{B_T} \rho(\xi)\int_{(Q_{\lceil r/\dej \rceil })_{\lambda_j}(\xi) \setminus (Q_{r/\dej})_{\lambda_j}(\xi)} f(x, x+\lambda_j\xi, M\xi)\, dx\,d\xi \\ \nonumber
    & \leq  \Bigl(\frac{\dej}{r}\Bigr)^d|E_j|\beta|M|^p\int_{B_T} \rho(\xi)|\xi|^p \, d\xi \\ \nonumber
    & =: \theta_j,
\end{align}
where we have set
\begin{equation}\label{set}
    E_j:=Q_{\lceil r/\dej \rceil } \setminus Q_{r/\dej-T\lambda_j}, \quad j\in\N,
\end{equation}
and we have used that $(Q_{\lceil r/\dej \rceil })_{\lambda_j}(\xi) \setminus (Q_{r/\dej})_{\lambda_j}(\xi)\subset E_j$ for all $\xi\in B_T$. Recalling \eqref{rho : moment} we infer there exists a positive constant $C$, independent of $j$, such that
\begin{equation*}
 \theta_j  = C \Bigl(\frac{\dej}{r}\Bigr)^d \Bigl( \Bigl\lceil \frac{r}{\dej} \Bigr\rceil^d - \Bigl(\frac{r}{\dej}-T\lambda_j\Bigr)^d \Bigr)  = C\Bigl[\Bigl( \frac{\lceil {r/\dej} \rceil}{r/\dej}\Bigr)^d - \Bigl(1-\frac{T}{r}\ej\Bigr)^d \Bigr],
\end{equation*}
which tends to $0$ as $j\to+\infty$. As a consequence of these observations, we get that
\begin{align*}
     &\Bigl(\frac{\dej}{r}\Bigr)^d\int_{B_T} \rho(\xi)\int_{(Q_{r/\dej})_{\lambda_j}(\xi)} f\Bigl(x,x+\lambda_j\xi,\frac{u(x+\lambda_j\xi)-u(x)}{\lambda_j}\Bigr)\, dx\,d\xi \\
     & \geq \Bigl(\frac{\dej}{r}\Bigr)^d\int_{B_T} \rho(\xi)\int_{(Q_{\lceil r/\dej \rceil })_{\lambda_j}(\xi)} f\Bigl(x,x+\lambda_j\xi,\frac{\widetilde{u}(x+\lambda_j\xi)-\widetilde{u}(x)}{\lambda_j}\Bigr)\, dx\,d\xi -\theta_j,    
\end{align*}
which, by the arbitrariness of $u$, yields
\begin{align*}
    \f(M) & = \lim_{j\to+\infty} \inf\Bigl\{ \Bigl(\frac{\dej}{r}\Bigr)^d\int_{B_T} \rho(\xi)\int_{(Q_{r/\dej})_{\lambda_j}(\xi)} f\Bigl(x,x+\lambda_j\xi,\frac{u(x+\lambda_j\xi)-u(x)}{\lambda_j}\Bigr)\, dx\,d\xi : \\
    & \qquad \qquad \qquad  u\in \mathcal{D}_{T\lambda_j,M}(Q_{r/\dej}; \R^m) \Bigr\} \\
    & \geq \limsup_{j\to+\infty} \inf\Bigl\{\Bigl(\frac{\dej}{r}\Bigr)^d\int_{B_T} \rho(\xi)\int_{(Q_{\lceil r/\dej\rceil})_{\lambda_j}(\xi)} f\Bigl(x,x+\lambda_j\xi,\frac{u(x+\lambda_j\xi)-u(x)}{\lambda_j}\Bigr)\, dx\,d\xi : \\
    & \qquad \qquad \qquad \quad u\in \mathcal{D}_{T\lambda_j,M}(Q_{\lceil r/\dej\rceil}; \R^m) \Bigr\}.
\end{align*}

Consider now $u\in \mathcal{D}_{T\lambda_j,M}(Q_{\lceil r/\dej\rceil};\R^m)$, let $\widetilde{w}$ be the $Q_{\lceil r/\dej\rceil}$-periodic extension of $u(x)-Mx$, let $\widetilde{u}(x):=\widetilde{w}(x)+Mx$, and define the function 
\begin{equation*}
    w(x):=\frac{1}{\lceil r/\dej \rceil^d} \sum_{i\in \mathbb{Z}^d\cap [0,\lceil r/\dej \rceil)^d} \widetilde{u}(x+i).
\end{equation*}
Since $w\in L^p_{\#,M}(Q_1; \R^m)$, by the assumptions \eqref{f : convex} and \eqref{f : period} we obtain
\begin{align} \nonumber
    \f_j(M) & \leq \int_{B_T} \rho(\xi)\int_{Q_1} f\Bigl(x,x+\lambda_j\xi,\frac{w(x+\lambda_j\xi)-w(x)}{\lambda_j}\Bigr)\, dx\,d\xi \\ \nonumber
    & \leq \frac{1}{\lceil r/\dej \rceil^d} \sum_{i\in \mathbb{Z}^d\cap [0,\lceil r/\dej \rceil)^d}  \int_{B_T} \rho(\xi)\int_{Q_1} f\Bigl(x, x+\lambda_j\xi, \frac{\widetilde{u}(x+i+\lambda_j\xi)-\widetilde{u}(x+i)}{\lambda_j}\Bigr)\, dx\,d\xi \\ \nonumber
    & = \frac{1}{\lceil r/\dej \rceil^d} \sum_{i\in \mathbb{Z}^d\cap [0,\lceil r/\dej \rceil)^d}  \int_{B_T} \rho(\xi)\int_{i+Q_1} f\Bigl(x, x+\lambda_j\xi, \frac{\widetilde{u}(x+\lambda_j\xi)-\widetilde{u}(x)}{\lambda_j}\Bigr)\, dx\,d\xi \\ \label{cell 1}
    & = \frac{1}{\lceil r/\dej \rceil^d} \int_{B_T} \rho(\xi)\int_{Q_{\lceil r/\dej\rceil}} f\Bigl(x,x+\lambda_j\xi,\frac{u(x+\lambda_j\xi)-u(x)}{\lambda_j}\Bigr)\, dx\,d\xi.
\end{align}
Moreover, observing that $Q_{\lceil r/\dej\rceil}\setminus (Q_{\lceil r/\dej\rceil})_{\lambda_j}(\xi)\subset E_j$ for all $\xi\in B_T$ and using \eqref{f : growth}, we get 
\begin{align} \nonumber
    & \frac{1}{\lceil r/\dej \rceil^d} \int_{B_T} \rho(\xi)\int_{Q_{\lceil r/\dej\rceil}\setminus (Q_{\lceil r/\dej\rceil})_{\lambda_j}(\xi)} f\Bigl(x, x+\lambda_j\xi,\frac{u(x+\lambda_j\xi)-u(x)}{\lambda_j}\Bigr)\, dx\,d\xi \\ \nonumber
    & \leq \frac{1}{\lceil r/\dej \rceil^d} |E_j| \beta |M|^p\int_{B_T} \rho(\xi)|\xi|^p\, d\xi \\ \label{cell 2}
    & = \frac{(r/\dej)^d}{\lceil r/\dej \rceil^d}  \theta_j.
 \end{align}
Therefore, combining \eqref{cell 1} with \eqref{cell 2}, we obtain
\begin{align*}
  \Bigl(\frac{\dej}{r}\Bigr)^d\int_{B_T} \rho(\xi)\int_{(Q_{\lceil r/\dej\rceil})_{\lambda_j}(\xi)} f\Bigl(x,x+\lambda_j\xi,\frac{u(x+\lambda_j\xi)-u(x)}{\lambda_j}\Bigr)\, dx\,d\xi \geq \Bigl(\frac{\lceil r/\dej \rceil}{r/\dej}\Bigr)^d\f_j(M)-\theta_j;
\end{align*}
hence, taking into account the arbitrariness of $u$ and that $\theta_j\to0$, we infer \eqref{cell lemma} passing to the limit, concluding the proof of $(ii)$.

\smallskip

Finally, we assume that  $\lambda\in[0,+\infty)$ and  we prove that
\begin{equation}\label{cell lemma bis}
      \f(M) \leq \liminf_{j\to+\infty} \f_j(M).
\end{equation}
Let $\{u_j\}_j\subset L^p_{\#,M}(Q_1;\R^m)$ be such that 
\begin{equation*}
    \liminf_{j\to+\infty} \f_j(M) = \liminf_{j\to+\infty} \int_{B_T} \rho(\xi)\int_{Q_1} f\Bigl(x,x+\lambda_j\xi,\frac{u_j(x+\lambda_j\xi)-u_j(x)}{\lambda_j}\Bigr)\, dx\,d\xi, 
\end{equation*}
and, by the translation invariance of the functional, suppose that
\begin{equation*}
    \int_{Q_1}u_j\,dx=0
\end{equation*}
for every $j\in\N$. Since $T>r_0$, by \eqref{rho : ball}, the lower bound in \eqref{f : growth}, and the previous step, we have 
\begin{equation*}
\f(M)\geq \alpha c_0\liminf_{j\to+\infty} \int_{B_{r_0}} \int_{Q_1}\Bigl|\frac{u_j(x+\lambda_j\xi)-u_j(x)}{\lambda_j}\Bigr|^p \, dx\,d\xi;
\end{equation*}
therefore, upon extracting a not relabeled subsequence,  we apply Lemma \ref{lemma : compactness} with $r=r_0$ and $A=Q_1$ to obtain that $\sup_{j} \|u_j\|_{L^p(Q_1;\R^m)}<+\infty$, and then
\begin{equation}\label{uniform bound}
 \sup_{j} \|u_j-Mx\|_{L^p(Q_1;\R^m)}<+\infty.
\end{equation}
Now we set $w_j(x):=\dej u_j(x/\dej), j\in \N$, and note that $w_j\to Mx$ in $L^p(Q_1;\R^m)$ as $j\to+\infty$. To see this, we also let $v_j(x):=u_j(x)-Mx, j\in\N$, and observe that, since $v_j$ is $Q_1$-periodic, \begin{align*}
    \int_{Q_1}|w_j(x)-Mx|^p\,dx & = \int_{Q_1}\Bigl|\dej v_j\Bigl(\frac{x}{\dej}\Bigr)\Bigr|^p\,dx \\
    & = \dej^{d+p}\int_{Q_{1/\dej}} |v_j(x)|^p\, dx \\
    & \leq \frac{\lceil1/\dej\rceil^d} {(1/\dej)^d}\dej^p\int_{Q_1}|v_j(x)|^p\,dx,
\end{align*}
which tends to $0$ by \eqref{uniform bound} and the fact that $\dej\to0$. Using the definition of $\Gamma$-limit and a change of variables, we obtain
\begin{align*}
  \f(M) = \int_{Q_1} \f(M)\, dx & \leq \liminf_{j\to+\infty} F^T_j(w_j, Q_1) \\
    & = \liminf_{j\to+\infty} \int_{B_T} \rho(\xi)\int_{{(Q_1)}_{\ej}(\xi)} f\Bigl(\frac{x}{\dej},\frac{x+\ej\xi}{\dej},\frac{w_j(x+\ej\xi)-w_j(x)}{\ej}\Bigr)\, dx\,d\xi \\
    & = \liminf_{j\to+\infty} \dej^d \int_{B_T} \rho(\xi)\int_{(Q_{1/\dej})_{\lambda_j}(\xi)} f\Bigl(x,x+\lambda_j\xi,\frac{u_j(x+\lambda_j\xi)-u_j(x)}{\lambda_j}\Bigr)\, dx\,d\xi \\
    & \leq \liminf_{j\to+\infty} \dej^d \int_{B_T} \rho(\xi)\int_{Q_{\lceil 1/\dej \rceil}} f\Bigl(x,x+\lambda_j\xi,\frac{u_j(x+\lambda_j\xi)-u_j(x)}{\lambda_j}\Bigr)\, dx\,d\xi \\
    & =  \liminf_{j\to+\infty} \int_{B_T} \rho(\xi)\int_{Q_1} f\Bigl(x,x+\lambda_j\xi,\frac{u_j(x+\lambda_j\xi)-u_j(x)}{\lambda_j}\Bigr)\, dx\,d\xi \\
    & = \liminf_{j\to+\infty}\f_j(M),
     \end{align*}
where we also used \eqref{f : period} and the $Q_1$-periodicity of the functions
\begin{equation*}
    x\mapsto \frac{u_j(x+\lambda_j\xi)-u_j(x)}{\lambda_j}, \quad j\in\N.
\end{equation*}
This leads to \eqref{cell lemma bis} and proves $(iii)$.
\end{proof}

\end{document}
